% GENERATED FILE. Edit canonical lessons, then run npm run generate:books.
% canonical-source-hash: fnv1a64:99d8c99ed2037a3b
% canonical-lessons: FA-C537-rang-kardan, FA-C537-kaghaz-divari, FA-C537-jabeja-kardan, FA-C537-kharab-shodan, FA-C537-sahebkhane

\chapter{Changes at Home}
\label{ch:fa-changes-at-home}
\hlchaptermodality{537}

\begin{chapteropening}
\textbf{By the end of this chapter:} \emph{I can talk about doing up a home.}

Complete the last lesson of chapter 537: I can talk about doing up a home.
\end{chapteropening}

\section[\texorpdfstring{\fa{رنگ} \fa{کردن} (rang kardan)}{رنگ کردن (rang kardan)}]{\fa{رنگ} \fa{کردن} (rang kardan) — to decorate with paint (a wall or room)}

\label{lesson:FA-C537-rang-kardan}



\begin{quote}
Before the new one: say the Persian for untidy, then the Persian for rubbish.
\end{quote}

\subsection*{You'll want to know: \fa{رنگ} \fa{کردن}}
\textbf{\fa{رنگ} \fa{کردن}} (\emph{rang kardan}) — \textquotedblleft{}to decorate with paint (a wall or room)\textquotedblright{}.

\begin{quote}
\fa{ما} \fa{دیوار} \fa{اتاق} \fa{را} \fa{آبی} \fa{رنگ} \fa{کردیم}.
\end{quote}

\begin{quote}
\emph{mâ divâr-e otâq râ âbi rang kardim.}
\end{quote}

\begin{quote}
We painted the wall of the room blue.
\end{quote}

\emph{rang} 'colour, paint' + \emph{kardan}. For pictures, \emph{naqqâshi kardan}.

\begin{cousinweb}
\emph{rang kardan} is built on two inherited words: \emph{rang}, 'colour', which EtymDB traces to a Middle Persian word for 'colour, dye' (EtymDB), and \emph{kardan}, 'to do'.
\end{cousinweb}

\begin{grammarlens}[title={What you've built}]
One more word for telling, explaining and describing at B1.
\end{grammarlens}

\subsection*{Your turn}
\begin{itemize}
\raggedright
  \item \emph{Say it:} \emph{rang kardan}
  \item \emph{Say it:} \emph{rang kardan}, once more
  \item \emph{Say it:} a sentence of your own with \emph{rang kardan}
  \item \emph{Recall:} say the Persian for untidy, then the Persian for rubbish, then say \emph{rang kardan} again
\end{itemize}

\begin{tcolorbox}[breakable,colback=teal!4,colframe=teal!35!black,title={Before you move on}]
What does \fa{رنگ} \fa{کردن} mean? (\textquotedblleft{}To decorate with paint (a wall or room)\textquotedblright{}.) Say it once more.
\end{tcolorbox}
\section[\texorpdfstring{\fa{کاغذ} \fa{دیواری} (kâghaz divâri)}{کاغذ دیواری (kâghaz divâri)}]{\fa{کاغذ} \fa{دیواری} (kâghaz divâri) — wallpaper}

\label{lesson:FA-C537-kaghaz-divari}



\begin{quote}
Before the new one: say the Persian for rubbish, then the Persian for to decorate with paint.
\end{quote}

\subsection*{You'll want to know: \fa{کاغذ} \fa{دیواری}}
\textbf{\fa{کاغذ} \fa{دیواری}} (\emph{kâghaz divâri}) — \textquotedblleft{}wallpaper\textquotedblright{}.

\begin{quote}
\fa{آنها} \fa{کاغذ} \fa{دیواری} \fa{اتاق} \fa{خواب} \fa{را} \fa{عوض} \fa{کردند}.
\end{quote}

\begin{quote}
\emph{ânhâ kâghaz divâri-ye otâq-e khâb râ 'avaz kardand.}
\end{quote}

\begin{quote}
They changed the wallpaper in the bedroom.
\end{quote}

\emph{kâghaz} 'paper' + \emph{divâri} 'of the wall'.

\begin{cousinweb}
\emph{kâghaz divâri} joins \emph{kâghaz}, 'paper', which EtymDB traces to Sogdian, another Iranian language (EtymDB), and \emph{divâri}, from \emph{divâr}, 'wall', which it traces to Middle Persian (EtymDB).
\end{cousinweb}

\begin{grammarlens}[title={What you've built}]
One more word for telling, explaining and describing at B1.
\end{grammarlens}

\subsection*{Your turn}
\begin{itemize}
\raggedright
  \item \emph{Say it:} \emph{kâghaz divâri}
  \item \emph{Say it:} \emph{kâghaz divâri}, once more
  \item \emph{Say it:} a sentence of your own with \emph{kâghaz divâri}
  \item \emph{Recall:} say the Persian for rubbish, then the Persian for to decorate with paint, then say \emph{kâghaz divâri} again
\end{itemize}

\begin{tcolorbox}[breakable,colback=teal!4,colframe=teal!35!black,title={Before you move on}]
What does \fa{کاغذ} \fa{دیواری} mean? (\textquotedblleft{}Wallpaper\textquotedblright{}.) Say it once more.
\end{tcolorbox}
\section[\texorpdfstring{\fa{جابجا} \fa{کردن} (jâbejâ kardan)}{جابجا کردن (jâbejâ kardan)}]{\fa{جابجا} \fa{کردن} (jâbejâ kardan) — to move (things), to rearrange}

\label{lesson:FA-C537-jabeja-kardan}



\begin{quote}
Before the new one: say the Persian for to decorate with paint, then the Persian for wallpaper.
\end{quote}

\subsection*{You'll want to know: \fa{جابجا} \fa{کردن}}
\textbf{\fa{جابجا} \fa{کردن}} (\emph{jâbejâ kardan}) — \textquotedblleft{}to move (things), to rearrange\textquotedblright{}.

\begin{quote}
\fa{ما} \fa{مبلها} \fa{را} \fa{جابجا} \fa{کردیم}.
\end{quote}

\begin{quote}
\emph{mâ moblhâ râ jâbejâ kardim.}
\end{quote}

\begin{quote}
We moved the sofas around.
\end{quote}

\emph{jâ} 'place' + \emph{be} 'to' + \emph{jâ} 'place': place to place. + \emph{kardan}.

\begin{cousinweb}
\emph{jâbejâ} is built on native Persian \emph{jâ}, 'place', which EtymDB traces to Middle Persian (EtymDB), said twice with \emph{be}, 'to': from place to place. \emph{kardan} is inherited.
\end{cousinweb}

\begin{grammarlens}[title={What you've built}]
One more word for telling, explaining and describing at B1.
\end{grammarlens}

\subsection*{Your turn}
\begin{itemize}
\raggedright
  \item \emph{Say it:} \emph{jâbejâ kardan}
  \item \emph{Say it:} \emph{jâbejâ kardan}, once more
  \item \emph{Say it:} a sentence of your own with \emph{jâbejâ kardan}
  \item \emph{Recall:} say the Persian for to decorate with paint, then the Persian for wallpaper, then say \emph{jâbejâ kardan} again
\end{itemize}

\begin{tcolorbox}[breakable,colback=teal!4,colframe=teal!35!black,title={Before you move on}]
What does \fa{جابجا} \fa{کردن} mean? (\textquotedblleft{}To move (things), to rearrange\textquotedblright{}.) Say it once more.
\end{tcolorbox}
\section[\texorpdfstring{\fa{خراب} \fa{شدن} (kharâb shodan)}{خراب شدن (kharâb shodan)}]{\fa{خراب} \fa{شدن} (kharâb shodan) — to break down; to go wrong}

\label{lesson:FA-C537-kharab-shodan}



\begin{quote}
Before the new one: say the Persian for wallpaper, then the Persian for to move.
\end{quote}

\subsection*{You'll want to know: \fa{خراب} \fa{شدن}}
\textbf{\fa{خراب} \fa{شدن}} (\emph{kharâb shodan}) — \textquotedblleft{}to break down; to go wrong\textquotedblright{}.

\begin{quote}
\fa{یخچال} \fa{دیروز} \fa{خراب} \fa{شد}.
\end{quote}

\begin{quote}
\emph{yakhchâl diruz kharâb shod.}
\end{quote}

\begin{quote}
The fridge broke down yesterday.
\end{quote}

\emph{kharâb} 'broken, ruined' + \emph{shodan} 'to become'.

\begin{cousinweb}
\emph{kharâb} comes from Arabic \emph{ḫarāb}, 'ruin' (EtymDB). \emph{shodan}, 'to become', is inherited Persian.
\end{cousinweb}

\begin{grammarlens}[title={What you've built}]
One more word for telling, explaining and describing at B1.
\end{grammarlens}

\subsection*{Your turn}
\begin{itemize}
\raggedright
  \item \emph{Say it:} \emph{kharâb shodan}
  \item \emph{Say it:} \emph{kharâb shodan}, once more
  \item \emph{Say it:} a sentence of your own with \emph{kharâb shodan}
  \item \emph{Recall:} say the Persian for wallpaper, then the Persian for to move, then say \emph{kharâb shodan} again
\end{itemize}

\begin{tcolorbox}[breakable,colback=teal!4,colframe=teal!35!black,title={Before you move on}]
What does \fa{خراب} \fa{شدن} mean? (\textquotedblleft{}To break down; to go wrong\textquotedblright{}.) Say it once more.
\end{tcolorbox}
\section[\texorpdfstring{\fa{صاحبخانه} (sâhebkhâne)}{صاحبخانه (sâhebkhâne)}]{\fa{صاحبخانه} (sâhebkhâne) — a landlord, a landlady; the host}

\label{lesson:FA-C537-sahebkhane}



\begin{quote}
Before the new one: say the Persian for to move, then the Persian for to break down.
\end{quote}

\subsection*{You'll want to know: \fa{صاحبخانه}}
\textbf{\fa{صاحبخانه}} (\emph{sâhebkhâne}) — \textquotedblleft{}a landlord, a landlady; the host\textquotedblright{}.

\begin{quote}
\fa{صاحبخانه} \fa{اجاره} \fa{را} \fa{بالا} \fa{برد}.
\end{quote}

\begin{quote}
\emph{sâhebkhâne ejâre râ bâlâ bord.}
\end{quote}

\begin{quote}
The landlord put the rent up.
\end{quote}

\emph{sâheb} 'owner' + \emph{khâne} 'house'. Also the host at a party.

\begin{cousinweb}
\emph{sâhebkhâne} is built in Persian on \emph{sâheb}, 'owner', from Arabic \emph{ṣāḥib}, 'companion' (EtymDB), and native \emph{khâne}, 'house'.
\end{cousinweb}

\begin{grammarlens}[title={What you've built}]
One more word for telling, explaining and describing at B1.
\end{grammarlens}

\subsection*{Your turn}
\begin{itemize}
\raggedright
  \item \emph{Say it:} \emph{sâhebkhâne}
  \item \emph{Say it:} \emph{sâhebkhâne}, once more
  \item \emph{Say it:} a sentence of your own with \emph{sâhebkhâne}
  \item \emph{Recall:} say the Persian for to move, then the Persian for to break down, then say \emph{sâhebkhâne} again
\end{itemize}

\begin{tcolorbox}[breakable,colback=teal!4,colframe=teal!35!black,title={Before you move on}]
What does \fa{صاحبخانه} mean? (\textquotedblleft{}A landlord, a landlady; the host\textquotedblright{}.) Say it once more.
\end{tcolorbox}
